\section{Conclusion}

The Paired Efficiency Factor $\eta=(1+\kappa)/(1+\kappa-2\sqrt{\kappa}\,\rho)$ generalises Fisher's classical paired efficiency to unequal-variance settings. The formula is distribution-free. Under bivariate normality, a functional relationship between $\eta$ and information content clarifies when and why relative features can improve predictions even when $\eta<1$: information content and efficiency measure distinct quantities that can diverge.

The primary validation, using KPI data from professional rugby union and football (\PEFtotalStudies{} studies across all four quadrants), yielded success rates of \PEFrugbySuccess--\PEFfootSuccess\%. Supporting validation across healthcare ($\eta=\PEFhealthMeanEta$), clinical genomics ($\eta=\PEFgenomicsMeanEta$), finance ($\eta=\PEFfinanceMeanEta$), and manufacturing ($\eta=\PEFmfgMeanEta$ on real {Bosch} {CNC} telemetry), summarised in \cref{tab:validation}, shows domain-specific heterogeneity in mean $\eta$ and in the proportion of units with $\eta>1$ (\PEFmfgSuccess--\PEFhealthSuccess\% in the supporting tier). The four-quadrant taxonomy provides feature engineering guidance, and the empirical PEF-to-ML mapping ($r=\PEFmlCorr$) allows estimation of expected performance gains.

Practices that appear domain-specific, market-adjusted returns, paired clinical designs, control charts, relative performance metrics, can be understood as instances of correlation-based variance reduction with different $(\kappa,\rho)$ values. The PEF framework connects classical statistical principles with modern machine learning practice for the question of when relative features outperform absolute features.
